\scp{Numerals}{08-04}
\citet[26]{Blust1981} sees the metathesis of PMP *isa ‘one’ in
both \ili{Soboyo} and \ili{Buru} as evidence of an exclusively shared innovation that
supports subgrouping them together.
\citet[35f]{Collins1981} adds to this the shared strategy of ‘minus-two
= 8’ and ‘minus-one = 9’ as further evidence in support of
subgrouping \ili{Buru}-Sula-\ili{Taliabu} together.
The fuller picture of numbers in \trf{08-14},
along with additional data already presented in \trf{07-14}
and \trf{07-15} reveal other considerations discussed below.


%\begin{table}\tcp{Selected numerals in \tsc{Taliabu} and nearby languages}{08-14}
%\begin{adjustbox}{width=\textwidth}
	%\begin{threeparttable}\st{2pt}
	%\begin{tabular}{llllllll}\lsptoprule
%*gloss	&	one												&	two								&	three					&	four										&	six								&	eight													&	nine\\
%PMP			&	*isa, *asa								&	*duha							&	*təlu					&	*əpat										&	*ənəm							&	*walu													&	*siwa\su{†}\\	\midrule
%\ili{Saluan}	&	{\hr}isaʔ									&	{\hr}ohua					&	{\hr}totoluʔ	&	{\hr}opat								&	{\hr}anom					&	{\hr}ualuʔ										&	{\hr}osioʔ\\
%\ili{Balantak}&	{\hr}isaʔ									&	{\hr}ruaʔ					&	{\hr}toluʔ		&	{\hr}paat								&	{\hr}noom		{\cdr}&	{\hr}waluʔ										&	{\hr}sioʔ\\
%\ili{Banggai}	&	\cg{(meeŋ)}								&	{\hr}lua					&	{\hr}tolu			&	\cg{(saŋkap)}						&	{\hr}noom		{\cdr}&	\hp{*w}alu										&	{\hr}sio\\
%\ili{Soboyo}	&	{\hr}sia						{\cdr}&	{\hr}howo		{\cdr}&	{\hr}tolu			&	{\hr}ŋ-haac				{\cdr}&	{\hr}nooŋ		{\cdr}&	{\hr}walu											&	(ta-sia){\cbl}\\
%\ili{Taliabu}	&	{\hr}sia						{\cdr}&	{\hr}howo		{\cdr}&	{\hr}tolu			&	{\hr}ŋ-haaʔ, ŋaa	{\cdr}&	{\hr}nooŋ		{\cdr}&	{\hr}walu											&	(ta-sia){\cbl}\\
%\ili{Kadai}		&	{\hr}sia						{\cdr}&	{\hr}howo		{\cdr}&	{\hr}tolu			&	{\hr}ŋaa					{\cdr}&	{\hr}noo		{\cdr}&	{\hr}walu											&	\cg{(ta-sia)}{\cbl}\\	\midrule
%\ili{Mangole}	&	{\hr}gia, ga-ia			{\cdr}&	{\hr}ga-ʔuu				&	{\hr}ga-telu	&	\cg{(ga-dia)}			{\cdr}&	{\hr}ga-nee	{\cdr}&	{\hr}ga-ta-hua	{\cbl}	&	(ga-ta-sia){\cbl}\\
%\ili{Sanana}	&	{\hr}hia						{\cdr}&	{\hr}ga-hu				&	{\hr}ga-tel		&	{\hr}ga-reha			{\cdr}&	{\hr}ga-nei	{\cdr}&	{\hr}ga-ta-hua	{\cbl}	&	(ga-ta-sia){\cbl}\\
%\ili{Lisela}	&	{\hr}m-sia-n, saa		{\cdr}&	{\hr}rua					&	{\hr}telo			&	{\hr}paa					{\cdr}&	{\hr}nee		{\cdr}&	{\hr}t-rua	{\cbl}			&	(ʧia /t-sia/){\cbl}\\
%\ili{Buru}		&	{\hr}em-sia-n, saa	{\cdr}&	{\hr}rua					&	{\hr}telo			&	{\hr}paa					{\cdr}&	{\hr}nee		{\cdr}&	{\hr}et-rua	{\cbl}			&	(eʧia /et-sia/) {\cbl}\\
%\ili{Ambelau}	&	{\hr}esbii								&	{\hr}lua					&	{\hr}relo			&	{\hr}faa					{\cdr}&	{\hr}nee		{\cdr}&	{\hr}βalo											&	{\hr}siβa\\	\midrule
%\ili{Kayeli}	&	{\hr}sileʔ {\tl} silei		&	{\hr}luaʔ					&	{\hr}telo			&	{\hr}haa					{\cdr}&	{\hr}nee		{\cdr}&	{\hr}waloʔ										&	{\hr}siwaʔ\\
		%\lspbottomrule
	%\end{tabular} 
		%\begin{tablenotes}\tnsize
			%\item[†]	{Despite the superficial similarity the
								%\tsc{Nuclear Sula-Buru}
								%words are not from *siwa but are phrases meaning
								%‘minus-one’ formed with reflexes
								%of *isa > **sia ‘one’.}
		%\end{tablenotes}
	%\end{threeparttable}
%\end{adjustbox}
%\end{table}

\begin{table}\tcp{Selected numerals in \tsc{Taliabu} and nearby languages}{08-14}
\begin{adjustbox}{width=\textwidth}
	\begin{threeparttable}\st{2pt}
	\begin{tabular}{lllllll}\lsptoprule
*gloss	&one											&	two							&	four									&	six								&	eight													&	nine\\
PMP			&*isa, *asa								&	*duha						&	*əpat									&	*ənəm							&	*walu													&	*siwa\su{†}\\	\midrule
\ili{Saluan}	&{\hr}isaʔ								&	{\hr}ohua				&	{\hr}opat							&	{\hr}anom					&	{\hr}ualuʔ										&	{\hr}osioʔ\\
\ili{Balantak}&{\hr}isaʔ								&	{\hr}ruaʔ				&	{\hr}paat							&	{\hr}noom		{\cdr}&	{\hr}waluʔ										&	{\hr}sioʔ\\
\ili{Banggai}	&\cg{(meeŋ)}							&	{\hr}lua				&	\cg{(saŋkap)}					&	{\hr}noom		{\cdr}&	\hp{*w}alu										&	{\hr}sio\\
\noalign{\vskip-1.5ex}
\multicolumn{7}{@{}c@{}}{\dashedline{127mm}} \\
\noalign{\vskip-0.5ex}
\ili{Soboyo}	&{\hr}sia						{\cdr}&	{\hr}howo	{\cdr}&	{\hr}ŋ-haac			{\cdr}&	{\hr}nooŋ		{\cdr}&	{\hr}walu											&	(ta-sia){\cbl}\\
\ili{Taliabu}	&{\hr}sia						{\cdr}&	{\hr}howo	{\cdr}&	{\hr}ŋ-haaʔ, ŋaa{\cdr}&	{\hr}nooŋ		{\cdr}&	{\hr}walu											&	(ta-sia){\cbl}\\
\ili{Kadai}		&{\hr}sia						{\cdr}&	{\hr}howo	{\cdr}&	{\hr}ŋaa				{\cdr}&	{\hr}noo		{\cdr}&	{\hr}walu											&	(ta-sia){\cbl}\\	\midrule
\ili{Mangole}	&{\hr}gia, ga-ia		{\cdr}&	{\hr}ga-ʔuu			&	\cg{(ga-dia)}		{\cdr}&	{\hr}ga-nee	{\cdr}&	{\hr}ga-ta-hua	{\cbl}	&	(ga-ta-sia){\cbl}\\
\ili{Sanana}	&{\hr}hia						{\cdr}&	{\hr}ga-hu			&	{\hr}ga-reha		{\cdr}&	{\hr}ga-nei	{\cdr}&	{\hr}ga-ta-hua	{\cbl}	&	(ga-ta-sia){\cbl}\\
\ili{Lisela}	&{\hr}m-sia-n, saa	{\cdr}&	{\hr}rua				&	{\hr}paa				{\cdr}&	{\hr}nee		{\cdr}&	{\hr}t-rua	{\cbl}			&	(ʧia /t-sia/){\cbl}\\
\ili{Buru}		&{\hr}em-sia-n, saa	{\cdr}&	{\hr}rua				&	{\hr}paa				{\cdr}&	{\hr}nee		{\cdr}&	{\hr}et-rua	{\cbl}			&	(eʧia /et-sia/) {\cbl}\\
\ili{Ambelau}	&{\hr}esbii								&	{\hr}lua				&	{\hr}faa				{\cdr}&	{\hr}nee		{\cdr}&	{\hr}βalo											&	{\hr}siβa\\	\midrule
\ili{Kayeli}	&{\hr}sileʔ {\tl} silei		&	{\hr}luaʔ				&	{\hr}haa				{\cdr}&	{\hr}nee		{\cdr}&	{\hr}waloʔ										&	{\hr}siwaʔ\\
		\lspbottomrule
	\end{tabular} 
		\begin{tablenotes}\tnsize
			\item[†]	{Despite the superficial similarity the
								\tsc{Nuclear Sula-Buru}
								words are not from *siwa but are phrases meaning
								‘minus-one’ formed with reflexes
								of *isa > **sia ‘one’.}
		\end{tablenotes}
	\end{threeparttable}
\end{adjustbox}
\end{table}

It was already shown in \srf{07-04-02},
(\trf{07-14} and \trf{07-15}) that numerals 1,
4, 6 are shaped CVV in \tsc{Sula-Buru},
and that this pattern is also found elsewhere.
We note that this pattern of metathesis to CVV in numerals 1,
4, 6, is also found in many Sulawesi languages (e.g.
\ili{Balantak} and \ili{Banggai}; \tsc{\ili{Bungku}-Tolaki} languages ‘four’
and ‘six’ [\citealp[156f]{Mead1999}]),
\tsc{Seram-Tanimbar-Bomberai} languages, and \tsc{Timor-Babar} languages.
We acknowledge that *isa > **sia is locally unique,
but it is also consistent with a widespread pattern of metathesis
to CVV in certain numerals.
For subgrouping purposes, we do not feel the lexically specific
metathesis of *isa > **sia outranks the strong differences in
the historical phonology of the behaviour of PMP *ə,
*-ay, *-aw, and other lexical
and morphological issues already discussed,
since numerals are a semantic domain that is easily borrowed,
and this lexically specific metathesis is more easily attributable
to known contact than the bundle of innovations in the historical
phonology discussed in \srf{08-01}.

Regarding numerals ‘eight’ and ‘nine’, we note four things.
First, the pattern of ‘minus-two = 8’
and ‘minus-one = 9’ does not cover the whole purported subgroup
of \tsc{Sula-Buru} + \tsc{Taliabu}.
The strategy of ‘minus-two = 8’ is found only in \tsc{Nuclear
Sula-Buru} languages, since reflexes of PMP *walu ‘eight’ are retained in all three
\tsc{Taliabu} languages, and \ili{Ambelau}.
A reflex of *siwa ‘nine’ is also retained in \ili{Ambelau}.
Secondly, the \tsc{Taliabu} node does not use this strategy for
both 8 and 9, since it preserves PMP *walu ‘eight’ in all three languages.
Thirdly, \ili{Mangole} and \ili{Sanana} both have \it{ga-ta-sia} with \it{s},
rather than the expected but unattested *[ga-ta-ia]
and *[ga-ta-hia] respectively, thus not reflecting the normal sound correspondences.
Fourthly, the strategy of ‘minus-two = 8’
and ‘minus-one = 9’ is also found elsewhere,
such as in many languages of South Sulawesi
\citep[130--133]{GrimesGrimes1987} with which the languages of this region have
been in contact, with such contact continuing into the present (see \frf{08-03}).
The uneven distribution of this strategy of ‘minus-two = 8’
and ‘minus-one = 9’ throughout the languages,
the uneven preservation of reflexes of PMP *walu ‘eight’
and *siwa ‘nine’, the irregular sound correspondences,
and the fact that the same strategy is found elsewhere,
make it more likely to be due to contact
and diffusion, rather than to shared inheritance from a common parent language.

\begin{figure}[h]\centering
	\frame{\includegraphics[width=\columnwidth]{08-03.jpg}}
	\fcp{Makassar \it{pinisi} being built on the north coast of \ili{Buru} 1984}{08-03}
\end{figure}
